\section{A smooth degeneration and its derivations}
\label{sec:degeneration}

We retain the notation of \Cref{sec:construction}, using the same letters
for elements of $P$ and their images in the domain $A=P/(H)$.  The
filtration constructed here will let us replace a locally nilpotent
derivation of $A$ by a homogeneous one on a smooth algebra with a simpler
equation.

\subsection{The valuation and its coefficients}

Consider the presented algebra
\begin{equation}\label{deg:S-R}
 S=\kk[p,x,y,z,u]/\bigl(xy-z(z+1)-p^3u\bigr),
 \qquad
 R=S/pS=\kk[x,y,z,u]/\bigl(xy-z(z+1)\bigr).
\end{equation}
The identity $xy-z(z+1)=p^2(H+pu)$ defines a homomorphism $S\to A$.

\begin{lemma}\label{deg:valuation}
The map $S\to A$ is injective, and
$S[p^{-1}]=A[p^{-1}]$.  The ring $R$ is a smooth domain.
On the common fraction field of $S$ and $A$ there is a discrete valuation
$\ord_p$, normalized by $\ord_p(p)=1$, coming from $S_{pS}$.
For $b\in S$ and every integer $i\geq 0$,
\begin{equation}\label{deg:global-divisibility}
 \ord_p(b)\geq i \quad\Longleftrightarrow\quad b\in p^iS.
\end{equation}
In particular, each nonzero $h\in S[p^{-1}]$ has a unique expression
$h=p^m b$ with $m=\ord_p(h)\in\mathbb Z$ and $b\in S\setminus pS$.
\end{lemma}

\begin{proof}
After inverting $p$, the defining polynomial of $S$ is linear in $u$
with invertible leading coefficient.  It is therefore irreducible in
$\kk[p,p^{-1},x,y,z,u]$.  It is not divisible by $p$ in
$\kk[p,x,y,z,u]$, so it is irreducible there as well: any factor becoming
a unit after localization would be a scalar times a power of $p$.
The latter polynomial ring is a unique factorization domain, so $S$ is
a domain and $p$ is nonzero in $S$.  The inverse formulas
\[
 s=y-x(x-u^3),\qquad
 F=\frac{x-s^2-u^3}{p^2},\qquad
 J=\frac{z-sx}{p^2}
\]
show that $S\to A$ becomes an isomorphism after inverting $p$.
Its kernel is consequently $p$-torsion, hence zero in the domain $S$.
This also proves equality of the fraction fields.

The polynomial $xy-z(z+1)$ is primitive as a polynomial in $y$ over
$\kk[x,z,u]$, since $x$ and $z(z+1)$ are coprime.  It is linear and
irreducible over $\kk(x,z,u)$, so $R$ is a domain.  Its hypersurface is
smooth by the Jacobian criterion~\cite[Tags~00T6 and~00T7]{Stacks}:
simultaneous vanishing of the partial
derivatives $y$, $x$, and $-(2z+1)$ would give $x=y=0$, $z=-1/2$,
which does not satisfy its equation.

Thus $pS$ is a nonzero prime ideal, and the noetherian local domain
$T=S_{pS}$ has maximal ideal $pT$.  To see directly that $T$ is a
discrete valuation ring, divide a nonzero nonunit repeatedly by $p$.
This process must end: otherwise the principal ideals generated by the
successive quotients form an ascending chain, and stabilization would
imply that $p$ is a unit, by cancellation in the domain $T$.
Every nonzero element of $T$ is therefore a power of $p$ times a unit.
This gives the stated valuation on $\Frac(T)$.

If $b\in S$ has positive order, then $b\in pT$, so $tb\in pS$ for
some $t\in S\setminus pS$.  Primality of $pS$ gives $b\in pS$.
Dividing by $p$ and repeating proves
\Cref{deg:global-divisibility}; its converse is immediate.
Finally, write $h=p^{-N}b_0$ with $b_0\in S$, $N\geq 0$, and apply
\Cref{deg:global-divisibility} to $b_0$ to divide out its full order.
The asserted expression follows, and its uniqueness follows from the
valuation and cancellation.
\end{proof}

For $0\ne h\in A$ define its \emph{valuation degree} by
$\deg h=-\ord_p(h)$, and put $\deg 0=-\infty$.  In particular,
\[
 \deg(hh')=\deg h+\deg h',\qquad
 \deg(h+h')\leq\max\{\deg h,\deg h'\}.
\]
These formulas give an exhaustive, separated filtration indexed by all
integers,
\[
 \mathcal F_dA=\{h\in A:\deg h\leq d\},\qquad
 \gr A=\bigoplus_{d\in\mathbb Z}\mathcal F_dA/\mathcal F_{d-1}A.
\]
Write $[h]_d$ for the class of $h\in\mathcal F_dA$ in the $d$th
summand, and $\inif(h)=[h]_{\deg h}$ for a nonzero $h$.
Give $R[\tau,\tau^{-1}]$ the grading in which $R$ has degree zero
and $\tau$ has degree $-1$.  \Cref{deg:valuation} gives a graded
injection
\begin{equation}\label{deg:graded-injection}
 \gr A\lhook\joinrel\longrightarrow R[\tau,\tau^{-1}],
 \qquad
 [h]_d\longmapsto \overline{p^dh}\,\tau^{-d}.
\end{equation}
Indeed $p^dh\in S$ whenever $h\in\mathcal F_dA$, including when
$d<0$, by the last assertion of \Cref{deg:valuation}.  Its residue
modulo $pS$ vanishes exactly when $h\in\mathcal F_{d-1}A$.
The displayed maps respect sums and products and are injective on
each summand.  This explains why their coefficients lie in $R$
itself.

\subsection{The exact associated graded algebra}

The initial forms of the five generators give a candidate for $\gr A$.
The point requiring proof is that no cancellation among their leading
terms produces an additional generator. We establish the presentation
of the candidate first, then reduce the weight of any representative
whose leading terms cancel.

In $R$ set
\begin{equation}\label{deg:center}
 s_R=y-x(x-u^3),\qquad
 f=x-s_R^2-u^3,\qquad
 g=z-s_Rx,\qquad I=(f,g).
\end{equation}

\begin{proposition}\label{prop:graded}
Under \Cref{deg:graded-injection}, the associated graded algebra is
the following specific subalgebra of $R[\tau,\tau^{-1}]$:
\begin{equation}\label{deg:G}
 G:=\gr A
   =R[\tau,I/\tau^2]
   =R[\tau,f/\tau^2,g/\tau^2].
\end{equation}
It is a smooth domain.  Give $P=\kk[p,s,u,F,J]$ the presentation
weights
\begin{equation}\label{deg:presentation-weights}
 (\omega(p),\omega(s),\omega(u),\omega(F),\omega(J))
       =(-1,0,0,2,2).
\end{equation}
Then every nonzero $h\in A$ has a polynomial representative whose
maximum weight is exactly $\deg h$, and there is a graded isomorphism
\begin{equation}\label{deg:graded-presentation}
 G\cong P/(H_{\mathrm{top}}),\qquad
 H_{\mathrm{top}}=x^2F-(1+2sx)J-p^2J^2,
 \quad x=s^2+u^3+p^2F.
\end{equation}
The images of $p,s,u,F,J$ in this isomorphism are respectively
$\tau,s_R,u,f/\tau^2,g/\tau^2$.
\end{proposition}

\begin{proof}
The residues of $s$ and $u$ in $R$ are the nonzero elements $s_R$
and $u$.  Also $f$ and $g$ are nonzero: specializing $x=z=0$ sends
$f$ to $-y^2-u^3$, and specializing $x=0$, $z=-1$ sends $g$ to
$-1$.  The identities $p^2F=x-s^2-u^3$ and $p^2J=z-sx$ therefore
give
\begin{equation}\label{deg:initial-generators}
 \begin{array}{c|ccccc}
 h&p&s&u&F&J\\ \hline
 \deg h&-1&0&0&2&2\\
 \inif(h)&\tau&s_R&u&f/\tau^2&g/\tau^2.
 \end{array}
\end{equation}
Here $s_R$ is nonzero because its specialization at $x=z=0$ is $y$.

Let $\theta:P\to R[\tau,\tau^{-1}]$ substitute the five initial
forms in \Cref{deg:initial-generators}.  Its image contains the
specified copy of $R$: the generators of that copy are recovered by
\begin{equation}\label{deg:recover-R}
 \begin{split}
 x&=s_R^2+u^3+\tau^2(f/\tau^2),\\
 y&=s_R+x(x-u^3),\\
 z&=s_Rx+\tau^2(g/\tau^2).
 \end{split}
\end{equation}
Consequently $\theta(P)$ is exactly the subalgebra on the right of
\Cref{deg:G}.  In particular this is an identification over the
given coefficient ring $R$.

For the presentation weights, $x$ has weight zero and
$H_{\mathrm{top}}$ is homogeneous of weight two, with
\begin{equation}\label{deg:top-relation}
 H=H_{\mathrm{top}}-pu,
 \qquad xy-z(z+1)=p^2H_{\mathrm{top}}.
\end{equation}
After inverting $p$, the coordinate change from
$(s,u,F,J)$ to $(x,y,z,u)$ and its inverse in
\Cref{deg:valuation} identify $\theta$ with the quotient map
\[
 \kk[p,p^{-1},x,y,z,u]\longrightarrow R[\tau,\tau^{-1}],
 \qquad p\longmapsto\tau.
\]
Its kernel is $(xy-z(z+1))$, or equivalently
$(H_{\mathrm{top}})$ in the original localized coordinates.
This is a prime ideal.  Moreover $p$ does not divide
$H_{\mathrm{top}}$, since its reduction modulo $p$ is
\[
 x_0^2F-(1+2sx_0)J,\qquad x_0=s^2+u^3,
\]
which is nonzero.  The same unique factorization argument as in
\Cref{deg:valuation} shows that $H_{\mathrm{top}}$ is irreducible,
hence prime, in $P$.  Its principal ideal is thus saturated with
respect to $p$: if $p^NQ\in(H_{\mathrm{top}})$ then
$Q\in(H_{\mathrm{top}})$.  Contracting the localized kernel yields
\begin{equation}\label{deg:kernel}
 \ker\theta=(H_{\mathrm{top}}).
\end{equation}

It remains to show that these initial forms generate all of $\gr A$.
For a nonzero polynomial $Q\in P$, let $\omega(Q)$ be its maximum
monomial weight.  If $Q$ represents a nonzero $h\in A$, then
$\deg h\leq\omega(Q)$ by \Cref{deg:initial-generators} and the
valuation inequalities.  Suppose the inequality is strict, and write
$M=\omega(Q)$.  The weight-$M$ part $Q_M$ has
$\theta(Q_M)=0$, since the degree-$M$ class of $h$ is zero.
By \Cref{deg:kernel},
\[
 Q_M=H_{\mathrm{top}}C,
\]
where $C$ may be taken homogeneous of weight $M-2$.  To justify the
homogeneity, decompose any quotient $C$ into its finitely many
homogeneous parts; multiplication by the nonzero homogeneous
$H_{\mathrm{top}}$ cannot cancel contributions of different weights
in the domain $P$.
The polynomial
\[
 Q-HC=Q-Q_M+puC
\]
represents the same element $h$ and has maximum weight strictly less
than $M$, because $puC$ has weight $M-3$.
Repeating this operation terminates: the maximum weights are strictly
decreasing integers bounded below by the fixed finite number $\deg h$.
Thus it ends with a representative of weight exactly $\deg h$.
No lower bound on the weights of all monomials in $P$ is needed.

The initial form of $h$ is now the image under $\theta$ of the top
part of such an exact representative.  This proves
$\gr A=\theta(P)$ and, together with \Cref{deg:kernel}, proves both
\Cref{deg:G} and \Cref{deg:graded-presentation}.

Finally, $G$ is a domain by its embedding in
$R[\tau,\tau^{-1}]$.  On the open set where $p$, or equivalently
$\tau$, is invertible, the presentation gives the smooth algebra
$R[\tau,\tau^{-1}]$.  On $p=0$, the partial derivatives of
$H_{\mathrm{top}}$ with respect to $F$ and $J$ reduce to
$x_0^2$ and $-(1+2sx_0)$.  They cannot vanish simultaneously, as
\[
 (1-2sx_0)(1+2sx_0)+4s^2x_0^2=1.
\]
The Jacobian criterion proves smoothness also at every point of
$p=0$, and hence on all of $\Spec G$.
\end{proof}

The weights in \Cref{deg:presentation-weights} compute the valuation
degree on $G$.  This degree will remain distinct from the fiber
weight and the auxiliary degree introduced later.

The description $G=R[\tau,I/\tau^2]$ is an affine modification in
the sense of Kaliman--Zaidenberg~\cite[Definition~1.1]{KalimanZaidenberg}:
the base ring is $R[\tau]$, the ideal is $(\tau^2,I)$, and the
element inverted in the construction is $\tau^2$. Here the exact
graded presentation, including its coefficient ring, is what permits
the later bundle pullback.

\subsection{Inducing locally nilpotent derivations}

We now transfer a derivation of $A$ to its associated graded algebra.
The exact representative statement in \Cref{prop:graded} is the needed
input: a bound on the five generators then bounds the degree shift on
every element of $A$.
This is the filtered-to-homogeneous derivation method developed by
Kaliman--Makar-Limanov~\cite[Sections~3--5]{KalimanMakarLimanov}.
We give the argument for the present integer-indexed filtration,
including the possibility of negative degrees.

A $\kk$-derivation of a $\kk$-algebra $C$ is a $\kk$-linear map
$D:C\to C$ satisfying $D(hh')=hD(h')+h'D(h)$.
It is \emph{locally nilpotent} if for each $h\in C$ there is an
integer $N\geq 1$ such that $D^Nh=0$; all subsequent iterates then
vanish as well.  We abbreviate this by LND.
On a graded algebra a derivation is \emph{homogeneous of degree
shift $j$} if it sends the component of degree $d$ to that of degree
$d+j$ for every $d\in\mathbb Z$.

\begin{lemma}\label{lem:induced-lnd}
Every nonzero LND $D$ of $A$ induces a nonzero homogeneous LND
$D_0$ of $G$.  If $0\ne h\in\ker D$, then
$\inif(h)\in\ker D_0$.
\end{lemma}

\begin{proof}
Let $(h_1,\ldots,h_5)=(p,s,u,F,J)$ be the generators of $A$ and set
\begin{equation}\label{deg:shift}
 j=\max_{Dh_i\ne0}\bigl(\deg(Dh_i)-\deg h_i\bigr).
\end{equation}
The maximum is over a finite nonempty set, since $D$ is nonzero.
We claim that
\begin{equation}\label{deg:shift-bound}
 \deg(Dh)\leq\deg h+j\qquad(0\ne h\in A).
\end{equation}
By \Cref{prop:graded}, choose a representative of $h$ of maximum
weight exactly $\deg h$.  Applying the product rule to each monomial,
every nonzero summand in its derivative has degree at most the
monomial's weight plus $j$: differentiating a factor $h_i$
changes its degree by at most $j$.  Summing these bounds gives
\Cref{deg:shift-bound}.

It follows that $D(\mathcal F_dA)\subseteq\mathcal F_{d+j}A$, so
\begin{equation}\label{deg:induced-map}
 D_0([h]_d)=[Dh]_{d+j}
\end{equation}
defines a homogeneous derivation of degree shift $j$ on $G$.
It is well defined because $D(\mathcal F_{d-1}A)$ lies in
$\mathcal F_{d+j-1}A$, and the product rule descends to the graded
algebra.  For a generator attaining the maximum in
\Cref{deg:shift}, the right side of \Cref{deg:induced-map} applied
to its initial form is nonzero.  Thus $D_0\ne0$.

For every $h\in\mathcal F_dA$ and $m\geq 0$, induction gives the
identity in the indicated graded component
\begin{equation}\label{deg:iterate}
 D_0^m([h]_d)=[D^mh]_{d+mj}.
\end{equation}
It includes the case when an intermediate class vanishes: the shift
bound ensures that a drop below the indicated degree persists in
all later iterates.  Local nilpotence of $D$ therefore kills each
homogeneous class after finitely many iterates of $D_0$.
Every element of $G$ is a finite sum of homogeneous classes, so a
maximum of their nilpotence bounds kills that element.  This proves
local nilpotence on all of $G$, without a sign restriction on $j$.
Finally, if $Dh=0$, \Cref{deg:induced-map} gives
$D_0\inif(h)=0$.
\end{proof}

The positive-degree-invariant obstruction is inspired by Gupta's
Derksen-invariant strategy; see \cite[Lemma~2.4 and the proof of
Proposition~3.6]{GuptaThreefold}. Here it is expressed directly through
polynomial coordinates. The characteristic-zero construction and the
subsequent bundle and rigidity arguments are proved in this paper;
Gupta's positive-characteristic theorem is not an input to them.

\begin{proposition}\label{prop:positive-invariant}
If $A\cong\kk\poly{4}$, then $G$ admits a nonzero homogeneous LND
with a nonzero homogeneous invariant of strictly positive valuation
degree.
\end{proposition}

\begin{proof}
Under the hypothesis choose polynomial coordinates
$q_1,q_2,q_3,q_4$ of $A$.
The subset $\mathcal F_0A$ is a $\kk$-subalgebra.  Since $F$ has
degree two by \Cref{deg:initial-generators}, this subalgebra is
proper.  Therefore at least one coordinate $q_r$ has positive
degree; otherwise the coordinates would generate $A$ inside
$\mathcal F_0A$.
Choose $i\ne r$.  The coordinate partial derivative
$D=\partial/\partial q_i$ is a nonzero LND of $A$ and satisfies
$Dq_r=0$.  By \Cref{lem:induced-lnd}, its induced derivation $D_0$
on $G$ is nonzero and homogeneous, and it kills the nonzero
homogeneous element $\inif(q_r)$ of degree $\deg q_r>0$.
\end{proof}

Neither \Cref{lem:induced-lnd} nor \Cref{prop:positive-invariant}
asserts that $D_0$ preserves the subalgebra $R\subset G$.
Coefficient preservation can fail even for a homogeneous LND on $G$.
The derivation $\Delta$ in \Cref{con:derivation} has presentation-weight
shift $-2$ and satisfies
$\Delta(H_{\mathrm{top}})=\Delta(H+pu)=p^3-p^3=0$.
It therefore descends through \Cref{deg:graded-presentation} to an LND
on $G$ sending $u\in R$ to $-\tau^2\notin R$.
This example concerns actions on $G$: $\Delta$ does not descend to
$A$, since $\Delta H=p^3$.
The next section lifts the action on $\Spec G$ directly, without
assuming $D_0(R)\subset R$.
